\section{讨论}
\label{sec:discussion}

\subsection{基线告诉了我们什么}
\label{sec:disc-insights}

实测结果支持五点方法学观察。

\paragraph{简单且设定良好的模型很强。}在 \task{2} 上，官方 Psych Sheet、Plackett--Luce 模型与核密度模拟在 seed-42 运行中给出完全相同的 Kendall $\kendall=0.7673$，三种子为 $0.8132\pm0.0342$。在 \task{3} 上，历史 \DNF{} 率（$0.3591$）与最佳树模型（$0.3731$）在统计上无法区分（$p=0.749$，效应可忽略），而贝叶斯 Beta--Binomial 后验（$0.3725$）在多种子均值上反而最好（$0.2938\pm0.0570$，树模型为 $0.2779\pm0.0724$）。两种情形下，简单且领域知识驱动的基线都能匹配、或几乎匹配更具表达力的替代方案。这并非放弃更复杂模型的理由，而是对期望的校准。只报告“模型 A 击败模型 B”的基准，会掩盖二者都可能在极小成本下被一个好先验匹配的事实。

\paragraph{通用深度与图模型在此表现不佳。}六大方法家族中有两个明显差于其更简单的对手，且多种子结果表明这一差距对评估集重采样是稳健的。在 \task{1} 上，LSTM（\maelog{} $=1.0440$，三种子 $0.9436\pm0.0754$）是\emph{最差}基线，甚至差于历史均值（$0.6657$），对树模型有大配对效应（$\Delta=+1.087$，$p=5.2\times10^{-8}$，Cohen's $d=0.90$）。在 \task{2} 上，GNN（$\kendall=0.5670$，三种子 $0.7440\pm0.1258$）落后 Psych Sheet $0.200$（$p=3.2\times10^{-7}$，$d=0.56$）。与此同时，规则感知特征结合树模型在 \task{1} 上胜出、在 \task{3} 上具竞争力。在此样本规模与特征预算下，领域结构——选手历史、轮次格式、项目身份——似乎比架构容量更重要，序列与图模型可能需要领域特定的归纳偏置才能具备竞争力。我们认为这是本基准最可操作的发现。

\paragraph{困难情形是结构性的，而非偶然的。}误差集中出现在数据稀疏或现象本身随机之处：\task{1} 的新手（\maelog{} $0.4151$）、\task{2} 的较晚切片（$\kendall=0.6437$）、\task{3} 的较晚切片与精英（\aucpr{} $0.2962$ 与 $0.1399$），以及各处冷启动选手（\maelog{} $0.9827$）。这些不是抽样偶然，而源于历史稀少、分布漂移与近乎确定性的执行。有用的基准因此必须依据硬子集来评判，这正是 \wcabench{} 将其设为一等公民的原因。

\paragraph{外推与识别是结论变脆弱之处。}\task{4} 与 \task{5} 正因难以验证而被纳入。\task{4} 的三阶估计从 $3.02$~秒（变点）、$3.29$~秒（GP+EVT）到 $3.74$~秒（指数衰减）不等，都远离社区锚点 $\approx 2.37$~秒。它们都不是统计意义上的“错误”——各自是其假设下的一致拟合——但合理方法之间的离散度本身就是发现，且三者的稳定性聚合在 $21$ 个项目上并不显著分离。类似地，识别出的迁移项目对从 $413$（相关）降至 $312$（因果森林）、$270$（IV）、$6$（DID；三种子均值 $2$）。报告分歧与不可识别，而非给出单一数字，才是科学上诚实的做法。

\paragraph{评估集方差不可忽视。}三种子重跑表明，有些结论稳健（树模型在所有种子下都以很大优势赢得 \task{1}），有些则不然（\task{3} 的冠军在树模型与 Beta--Binomial 先验之间翻转，GNN 在 \task{2} 上的离散度 $\pm0.1258$ 大于若干模型间差距）。因此，本领域的单次排行榜应结合多种子表一同解读。

\subsection{关于避免基准饱和}
\label{sec:disc-saturation}

基准会退化：一旦任务被解决，它就不再具有区分度。我们设置了两道防线。第一，五个任务构成从设定良好的回归到因果识别的难度谱，因此易端的提升不会耗尽基准。第二，每个任务都附带硬子集（第~\ref{sec:protocol-stratification} 节）；对 \task{3}，限定历史 \DNF{} 率位于 $[0.1,0.3]$ 可去除易分类的极端；对 \task{2}，近似平局子集去除了排序平凡的轮次。小样本结果已表明若干子集远未饱和，例如精英 \task{3}（\aucpr{} $0.1399$）与冷启动 \task{1}（\maelog{} $0.9827$），因此基准仍有余量。

\subsection{为何需要一个领域基准，为何选速拧}
\label{sec:disc-domain}

强简单模型与更具表达力的替代方案并存，令人想起数据建模与算法预测之间的长期张力 \citep{breiman_cultures}：基准的职责是让比较公平，而非钦定胜者。\wcabench{} 是一个更宏大理念的案例研究：\emph{受规则约束的个人项目}是机器学习尚未充分利用的试验场。此类数据兼具规模与真实结构——显式格式、确定性平均规则、离散状态码——这些对通用管线不可见，却对准确性具有决定性。速拧是一个格外干净的实例：完全个人化、精确到百分之一秒计时、全球分布且公开发布。相关经验可迁移到游泳、田径或力量举等具有类似轮次格式与取消资格规则的领域。因此我们把 \wcabench{} 视为起点而非终点，视为一个在受规则约束的竞技领域中构建可信基准的可复用模板。

\subsection{开放问题}
\label{sec:disc-open}

若干方向被刻意留待后续。序列模型（LSTM/Transformer）与图神经网络目前表现不佳，因此开放问题不是它们原则上是否更准，而是哪些领域特定归纳偏置——项目感知嵌入、轮次格式约束、以正面对抗历史为锚的交互图——能让它们利用树模型已经捕捉到的结构。在 \task{5} 上，从 $413$ 个相关对到 $270$ 个 IV 对、再到 $6$ 个 DID 对的陡降，促使我们寻求更可信的因果设计（例如以某项目在各国错峰首次举办作为工具变量，或合成控制设计），而验证它们的困难又促使更好的敏感性分析。最后，非平稳性与项目尺度的相互作用表明，按项目的分层归一化（而非单一全局预处理）可能是最有前景的通用改进。
